\section{Maximum-depth scheduling ingredients}
\label{sec:worst}

The maximum-depth algorithm of Section~\ref{sec:polynomial-slots}
constructs its pattern children first and then schedules their discovery.
This section isolates its local tools. The priority scheduler extends
the construction beyond the exact scheduler's logarithmic rank cap by
supplying a passing height trial. All branch decisions use hypothesis
counts and supplied-line geometry.

Let $\delta\in(0,1/3)$ be the count-balance threshold and let $\gamma>2$
be the coefficient of the count potential
$\Phi(s)=\lfloor\gamma\log_2s\rfloor$. Assume
\begin{equation}\label{eq:worst-parameter-conditions}
 \gamma\log_2\frac1{1-\delta}\ge1,\qquad
 \gamma\log_2\frac1\delta\ge5.
\end{equation}
The first inequality pays for a balanced query; the second pays for
endpoint recovery at a flat state. Section~\ref{sec:worst-calibration}
chooses the constants at the point where these two credits balance.

For a state $S$, write $C_W(S)$ for the height of the tree returned
by the hybrid recursive algorithm. When a proof refers to a stored
child code, its height is the actual height of that code.

\subsection{Local credits and the priority fallback}

\begin{proofstrategy}
\begin{enumerate}
\item Restrict an optimal tree to a pattern class: its common witness
prefix disappears and gives that class a discovery deadline.
\item Recover the pattern with supplied endpoint comparisons; the
count-potential drop pays for this decoder.
\item Use priorities to meet all deadline prefixes in one ordering.
Every accepted ordering is checked against actual stored child heights.
\end{enumerate}
\end{proofstrategy}

In a count-flat state $S$ of size $s$, use the common witness and hull
patterns of Lemmas~\ref{lem:helly} and~\ref{lem:hull-structure}.
Every nonzero class $A$ has $a=|A|<\delta s$; the zero class $A_0$
is empty or a singleton. Write
\[
 z_A=\lfloor\log_2(\delta s/a)\rfloor.
\]
All child heights below come from already constructed codes; the
optimal tree appears only in the comparison argument.

The recursion needs integer query allowances, while its potential uses
logarithms. The following credit converts size reduction into a whole
query at a balanced split and enough whole queries for a flat decoder.

\begin{proposition}[Integer count credit]\label{lem:worst-floor-credit}
A balanced child of size $a\le(1-\delta)s$ satisfies
$\Phi(s)-\Phi(a)\ge1$. A nonzero flat-state class satisfies
\begin{equation}\label{eq:worst-floor-credit}
 \Phi(s)-\Phi(a)\ge\lfloor\gamma\log_2(s/a)\rfloor\ge z_A+5.
\end{equation}
\end{proposition}

For a flat state, the optimal tree supplies an analytical comparison
schedule. Its initial witness answers are constant on a class until
that class meets a support; restriction turns this delay into saved
height inside the class.

\begin{proposition}[Supporting-prefix credit]\label{lem:height-prefix}
Follow a common witness in a proper optimal tree of height
$D=\OPT_W(S)$ along a path of length $r$. Let $t_A$ be the first
path query whose real hull supports meet the positive pattern of $A$;
put $t_A=r+1$ if none does. Then
\[
 \OPT_W(A)\le D-(t_A-1).
\]
The supporting singleton or adjacent-pair blocks cover every nonzero
class by its index $t_A$, after appending one positive support for an
uncovered terminal nonzero singleton when necessary.
\end{proposition}

A first positive answer locates one vertex of the pattern. The endpoint
decoder then recovers the full interval, with a pointwise cost bound so
maximum depth controls each class separately.

\begin{proposition}[Legal endpoint allowances]\label{lem:endpoint-height}
After a first positive hull answer, the pointwise alphabetic endpoint
codes recover class $A$ with at most
\[
 e_A^{(4)}=z_A+4
\]
additional supplied queries. If a neighboring hull vertex is already
known negative, one endpoint is fixed and the allowance is
$e_A^{(2)}=z_A+2$. A hull of dimension at most one needs no further
query after discovery.
\end{proposition}

Scheduling can now use the codes already built for the children.
Adding their actual heights to discovery and endpoint recovery gives
a directly checkable route certificate, including the all-negative route.

\begin{proposition}[Local route-height certificate]\label{lem:worst-local-height}
Let $h_A$ be a stored class code's actual height. Discover a positive
vertex in a fixed ordering truncated when all nonzero patterns are
covered. If its first positive position for $A$ is $g_A$, the whole
route through $A$ has height at most $g_A+h_A+e_A^{(4)}$, or
$g_A+h_A+e_A^{(2)}$ on a pinned exit. If every such bound is at most
$H$, the whole tree has height at most $H$, including its zero route
when that class exists.
\end{proposition}

The classes have different deadlines, yet the algorithm must choose
one query order. Integer priorities make later deadlines unable to
interfere with greedy progress on any earlier deadline prefix.

\begin{proposition}[One ordering for all block deadlines]
\label{lem:deadline-cover}
Force an oriented adjacent real pair $(v,w)$ first. Consider only
classes negative on both, with integer block deadlines $d_A\ge2$.
Suppose a comparison sequence begins with $\{v,w\}$ and hits $A$
within $d_A$ singleton or adjacent-pair blocks. Sort the distinct
budgets as $d^{(1)}<\cdots<d^{(r)}$, give a deadline-$d^{(j)}$ class
weight $(s+1)^{r-j}$, and complete by maximum-remaining-weight greedy.
This polynomial-time ordering satisfies $g_A\le3d_A-2$
for every surviving class simultaneously.
\end{proposition}

The endpoint and prefix credits turn the optimal witness path into
deadlines, and the priority ordering meets those deadlines. It therefore
provides a passing trial even when the capped exact scheduler cannot
certify the stronger bound.

\begin{proposition}[Priority fallback height]
\label{thm:worst}
Suppose stored pattern children satisfy the fallback bound
\[
h_A\le3\OPT_W(A)+\Phi(|A|).
\]
At a flat state,
Algorithm~\ref{alg:priority-trial} has a passing trial at
\[
 H^*=3\OPT_W(S)+\Phi(s).
\]
A balanced query also fits this height if its two stored children
obey the same bound. Consequently the hybrid recursion can always
retain $C_W(S)\le3\OPT_W(S)+\Phi(s)$ as its fallback invariant.
\end{proposition}

For a trial $H$ and a forced pair, first-positive exits must pass
$1+h_A+e_A^{(4)}\le H$, and second-query exits must pass
$2+h_A+e_A^{(2)}\le H$. Classes negative on both receive
\[
 u_A=h_A+e_A^{(4)},\qquad B_A=H-u_A,\qquad
 d_A=\lfloor(B_A+2)/3\rfloor.
\]
Reject any $B_A<4$ and accept only after checking the resulting
actual first-hit positions against $g_A+u_A\le H$.
These tests describe an executable height certificate; the optimal
comparison blocks are used only to prove that some trial passes.

For a fixed target height $H$, each oriented adjacent pair determines
which patterns are hit at its first endpoint, second endpoint, or
neither. The first two groups have immediate local height tests, so a
pair that already exceeds $H$ can be discarded. For patterns left on
the miss route, the remaining height becomes a deadline; the priority
weights make the greedy order favor the tightest deadlines. The final
check uses the actual hit ranks and child allowances, and accepts the
order only when every miss-route pattern fits.

\begin{algorithm}[H]
\caption{\PriorityTrial: internal fallback for one height trial}
\label{alg:priority-trial}
\KwIn{Flat state's hull and classes; child heights $h_A$; allowances $e_A^{(4)},e_A^{(2)}$; target $H$.}
\KwOut{A passing discovery order, or infeasible.}
\ForEach{oriented adjacent real pair $(v,w)$}{
  Form $\mathcal F,\mathcal S,\mathcal N$ by their first hit on $(v,w)$\;
  \If{some $A\in\mathcal F$ has $1+h_A+e_A^{(4)}>H$, or
  some $A\in\mathcal S$ has $2+h_A+e_A^{(2)}>H$}{\Continue\;}
  \ForEach{$A\in\mathcal N$}{
    $u_A\gets h_A+e_A^{(4)}$;
    $B_A\gets H-u_A$;
    $d_A\gets\lfloor(B_A+2)/3\rfloor$\;
  }
  \If{some $B_A<4$}{\Continue\;}
  Sort distinct budgets as $d^{(1)}<\cdots<d^{(r)}$\;
  Give each deadline-$d^{(j)}$ class weight $(|S|+1)^{r-j}$\;
  Force $v,w$ and complete by maximum-remaining-weight greedy;
  truncate when all nonzero classes are hit\;
  \If{every $A\in\mathcal N$ has actual $g_A+u_A\le H$}{\KwRet this order\;}
}
\KwRet infeasible\;
\end{algorithm}


\subsection{Technical proofs of the local ingredients}

\begin{proof}[Proof of Proposition~\ref{lem:worst-floor-credit}]
For real $x,y$, write $x-y=k+\theta$ with integer
$k=\lfloor x-y\rfloor$ and $0\le\theta<1$. Then
$\lfloor x\rfloor=\lfloor y+k+\theta\rfloor
\ge\lfloor y\rfloor+k$, proving the floor-difference inequality.
Apply it to the two count potentials:
\[
 \Phi(s)-\Phi(a)
 \ge\left\lfloor\gamma\log_2(s/a)\right\rfloor.
\]
Indeed, the difference of the unfloored logarithms is
$\gamma\log_2s-\gamma\log_2a=\gamma\log_2(s/a)$, which is the
argument to the floor-difference inequality. For a balanced child,
$a\le(1-\delta)s$, hence $s/a\ge1/(1-\delta)$. The first
condition in \eqref{eq:worst-parameter-conditions} therefore gives
\[
 \gamma\log_2(s/a)
 \ge\gamma\log_2\frac1{1-\delta}\ge1.
\]
Since every real number at least one has floor at least one, the
balanced-child claim follows, including the equality case.
For a flat class put $x=\log_2(s/a)$ and
 $\ell=\log_2(1/\delta)$. Since $a<\delta s$, we have $x>\ell$.
The calibration gives $\gamma\ell=5$, so, using $\gamma>1$,
\[
 \gamma x=\gamma\ell+\gamma(x-\ell)
          =5+\gamma(x-\ell)>5+(x-\ell).
\]
Monotonicity of the floor therefore gives
\[
 \left\lfloor\gamma x\right\rfloor
 \ge\lfloor5+x-\ell\rfloor
 =5+\lfloor\log_2(\delta s/a)\rfloor=5+z_A,
\]
which proves the flat credit including its integer rounding.
\end{proof}

\begin{proof}[Proof of Proposition~\ref{lem:height-prefix}]
There are two distinct steps. Common answers give a restriction credit;
supporting blocks then cover the positive classes with deadlines measured
from the same witness path. The terminal leaf needs separate attention
because the witness may end at a real hypothesis with a nonzero pattern.
Claims~\ref{clm:height-prefix-constant} and
\ref{clm:height-comparison-cover} establish these two steps, respectively.

\begin{claim}[The common prefix disappears under restriction]
\label{clm:height-prefix-constant}
Every member of $A$ follows the witness's first $t_A-1$ answers.
\end{claim}
\begin{proof}
For every index $i<t_A$, no support of the $i$th query meets the positive
pattern of $A$. Lemma~\ref{lem:hull-structure} then makes its answer
constant on $A$, with the same value as on the witness. This includes a
query whose support uses the free origin marker. Thus every member of $A$
traverses the same first $t_A-1$ nodes of the optimal tree.

Restrict that tree to $A$. Since $A$ is nonempty and each of those answers
is constant, each corresponding node contracts to its witness-side child.
The remaining tree identifies $A$ and has height at most
$D-(t_A-1)$: each original path had length at most $D$, and the shared
prefix has been removed from every such path. Minimizing over identifying
trees on $A$ proves the inequality.

\end{proof}
The restriction credit must come with a cover that can be compared
against a scheduler. Queries along the path cover every class except,
possibly, the hypothesis at its terminal leaf.

\begin{claim}[The supporting blocks form a comparison cover]
\label{clm:height-comparison-cover}
Only an uncovered terminal nonzero singleton needs an appended block.
\end{claim}
\begin{proof}
Each class with $t_A\le r$ is hit by that path query's supports.
By Claim~\ref{clm:height-prefix-constant}, an unhit class follows
the entire original path and is therefore the
single leaf hypothesis reached by the witness. If its pattern is
nonzero, append a positive hull vertex at index $r+1$; if it is zero,
the witness leaf already completes the cover. Retaining the first
occurrence of a repeated support preserves every first-hit deadline.
\end{proof}
\end{proof}

\begin{proof}[Proof of Proposition~\ref{lem:endpoint-height}]
The query-free cut of Lemma~\ref{lem:legal-endpoint-queries} makes the
positive pattern an ordinary interval. Let $R$ be the hypotheses
compatible with the first positive transcript, so $|R|<\delta s$.
Use uniform counts, conditioned on that transcript, as code-design
probabilities, while the algorithm itself follows the deterministic
construction.
The two conditional endpoint probabilities multiply to the probability
of the full class. Their logarithmic query bounds therefore telescope
to one bound depending on the class size.
Claims~\ref{clm:height-two-endpoints} and~\ref{clm:height-pinned-endpoint}
establish the two allowances, according to whether discovery has
already fixed a boundary.

\begin{claim}[Two endpoint searches telescope their probabilities]
\label{clm:height-two-endpoints}
If $a_L$ compatible hypotheses have the target's left endpoint, its
full pattern is learned in fewer than $\log_2(\delta s/a)+4$ queries.
\end{claim}
\begin{proof}
Lemma~\ref{lem:legal-endpoint-queries} makes both searches legal.
Applying the pointwise bound of Lemma~\ref{lem:alphabetic-endpoints}
to the first endpoint and then to the conditional second endpoint gives
\[
 \log_2\frac{|R|}{a_L}+2+
 \log_2\frac{a_L}{a}+2
 =\log_2\frac{|R|}{a}+4<\log_2\frac{\delta s}{a}+4.
\]
The logarithms add because
$(|R|/a_L)(a_L/a)=|R|/a$: the second distribution is conditioned
on the learned first endpoint. The endpoint pair uniquely specifies
the pattern, while the strict inequality uses $|R|<\delta s$:
the first positive transcript confines $R$ to one minority cap.
If $q_A$ is the integer number of endpoint queries,
the displayed strict bound gives
$q_A<\log_2(\delta s/a)+4<z_A+5$. Thus
$q_A\le z_A+4$, including when the logarithm itself is an integer.
\end{proof}
A negative neighbor supplies more information than ordinary discovery:
it fixes an endpoint. This removes an entire endpoint search and gives
the extra saving needed on the second exit of the initial pair.

\begin{claim}[A cached negative neighbor removes one search]
\label{clm:height-pinned-endpoint}
A pinned positive vertex fixes one interval endpoint.
\end{claim}
\begin{proof}
Cut at the negative neighbor and orient the order to place the
positive vertex first. Lemma~\ref{lem:legal-endpoint-queries}
fixes that boundary; hence Lemma~\ref{lem:alphabetic-endpoints}
requires only one search to learn the other endpoint
in fewer than $\log_2(\delta s/a)+2$ queries. Writing its integer
query count as $q_A$, we have
$q_A<\log_2(\delta s/a)+2<z_A+3$, hence $q_A\le z_A+2$.
The negative neighbor is a real cached answer, so no anchor
query is added. In dimension at most one the positive pattern is
already determined; inference and cached answers can only lower cost.
\end{proof}
\end{proof}

\begin{proof}[Proof of Proposition~\ref{lem:worst-local-height}]
A route through $A$ makes at most $g_A$ discovery queries,
then uses the endpoint allowance of Proposition~\ref{lem:endpoint-height},
and finally follows a stored
code of height $h_A$. Its height is their sum, bounded by $H$ by
hypothesis. The ordering is truncated as soon as all nonzero classes
are covered, namely after at most $\max_{A\ne A_0}g_A$ queries.
Since $h_A$ and both endpoint allowances are nonnegative, the route
tests imply $g_A\le H$ for each nonzero class. Thus the all-negative
route also has length at most $H$ and ends at the zero singleton. If that class is absent its leaf is
omitted. A nonsingleton state has a nonzero pattern because
Lemma~\ref{lem:hull-structure} gives $|A_0|\le1$.
\end{proof}

\begin{proof}[Proof of Proposition~\ref{lem:deadline-cover}]
The proof separates priorities from geometry. First, the priorities
make the global order greedy for every nested deadline prefix. Then
an interval-cover bound limits the number of queries for each prefix.

\begin{claim}[Later deadlines cannot change a prefix greedy choice]
\label{clm:deadline-priority-prefix}
The global ordering is weighted greedy for every surviving deadline
prefix until that prefix is covered.
\end{claim}
\begin{proof}
Put $U=(s+1)^{r-j}$. A class of earlier or equal deadline has
weight $U(s+1)^{j-i}$ for some $i\le j$, so every candidate's
prefix score is an integer multiple of $U$. For $j<r$, all later
classes together have weight at most
\[
 s(s+1)^{r-j-1}=\frac{s}{s+1}U<U;
\]
for $j=r$ their total is zero. Deleting already hit classes can only
lower these totals. If one candidate's prefix score exceeds another's,
it exceeds it by at least $U$; no difference in their later-class
scores can reverse that comparison. Therefore a global score maximizer
also maximizes the prefix score. A nonzero prefix score is at least
$U$ and beats every zero prefix score while demands remain. Global
tie-breaking chooses among legitimate prefix greedy ties.
This applies to every deadline prefix at every step of the same order.
\end{proof}
The prefix argument determines which demands each greedy choice serves.
To bound the number of choices, use the forced adjacent pair to cut the
line of blocks and remove its first block; the remaining cover's size
and runs then pay for every later greedy choice.

\begin{claim}[The first pair leaves at most $3d-4$ later greedy choices]
\label{clm:deadline-prefix-length}
A nonempty deadline-$d$ prefix finishes by actual position $3d-2$.
\end{claim}
\begin{proof}
Cut at $v$ and delete its adjacent vertex $w$. Surviving demands are
ordinary intervals. Remove these two vertices from the union $P$
of the first $d$ comparison blocks. Each of the other $d-1$ blocks
contributes at most three to $|P|+c(P)$, where $c(P)$ counts runs:
a pair contributes at most two vertices and one run, and a singleton
at most one vertex and one run. Deleting the two adjacent cut vertices
cannot split any surviving adjacent-pair block. For finite sets
$U,V$ on a line, $|U\cup V|\le|U|+|V|$ and
$c(U\cup V)\le c(U)+c(V)$: joining runs cannot create a new run.
Applying these inequalities to the remaining $d-1$ blocks gives
$|P|+c(P)\le3(d-1)$. Every prefix demand misses the forced pair
and was hit by one of these remaining blocks, so $P$ still covers it.
By Claim~\ref{clm:deadline-priority-prefix}, the actual later
choices form a greedy sequence for this prefix. Consequently
Lemma~\ref{lem:greedy-cover-length} bounds their number by
\[
 |P|+c(P)-1\le3(d-1)-1=3d-4.
\]
The inequality uses the size-plus-run bound just established;
the equality simplifies the resulting integer allowance.
Adding the two forced queries gives position at most $3d-2$. An empty prefix needs no further choice.
\end{proof}
The auxiliary weights have $O(s\log s)$ bits, and their incidence
scores are computed from the explicit class-to-vertex table. Thus the
ordering and actual first-hit times have polynomial bit complexity.
\end{proof}

\begin{proof}[Proof of Proposition~\ref{thm:worst}]
Put $D=\OPT_W(S)$. For a balanced child $A$, restriction and
count credit, specifically Lemma~\ref{lem:restriction} and
Proposition~\ref{lem:worst-floor-credit}, give
\[
 \OPT_W(A)\le D,\qquad \Phi(|A|)\le\Phi(s)-1.
\]
Thus its complete route satisfies
\[
 1+h_A\le1+3\OPT_W(A)+\Phi(|A|)
          \le3D+\Phi(s)=H^*.
\]
The first inequality is the stored-child invariant. The second
substitutes the two preceding bounds, so the potential drop pays
for the new root query.
At a flat state,
enumerate an initial pair containing the first optimal-witness query's
supports; pad a singleton support by an adjacent real neighbor.
The first two exits must be checked directly: the deadline scheduler
handles only classes that miss both initial queries. The cached
negative answer makes the second positive exit cheaper to decode.

\begin{claim}[Both initial positive exits fit $H^*$]
\label{clm:fallback-initial-exits}
The first-query exit has at least one query of discovery credit;
the pinned second-query exit has at least three.
\end{claim}
\begin{proof}
For either class, Lemma~\ref{lem:restriction} gives
$\OPT_W(A)\le D$. Using the stored-child invariant and then
Proposition~\ref{lem:worst-floor-credit},
\[
 H^*-h_A\ge3D+\Phi(s)
                -3\OPT_W(A)-\Phi(|A|)
             \ge\Phi(s)-\Phi(|A|)
             \byref{Prop.}{lem:worst-floor-credit}{\ge}z_A+5.
\]
The middle inequality drops the nonnegative difference
$3(D-\OPT_W(A))$. Thus the first positive exit satisfies
$1+h_A+(z_A+4)\le H^*$. By Proposition~\ref{lem:endpoint-height},
the second positive exit has a cached negative neighbor and only
the smaller allowance, so
$2+h_A+(z_A+2)\le H^*-1\le H^*$.
Neither calculation assumes that the class meets the original support;
restriction and flat count credit also cover classes hit only by
the padding vertex.
\end{proof}
For the remaining classes, the supporting-prefix credit becomes an
integer block deadline. The floor below preserves that credit exactly,
so the comparison schedule meets the deadlines computed from actual
child heights.

\begin{claim}[Every later class meets its computed deadline]
\label{clm:fallback-later-deadline}
For a class negative on the initial pair, $d_A\ge t_A$.
\end{claim}
\begin{proof}
Here $t_A\ge2$, because this class misses the first query's
supports. Proposition~\ref{lem:height-prefix} and the actual
child invariant give
$h_A\le3[D-(t_A-1)]+\Phi(|A|)$. Substituting this into the
budget and then using count credit yields
\begin{align*}
 B_A&=3D+\Phi(s)-h_A-(z_A+4)\\
    &\ge3(t_A-1)+\Phi(s)-\Phi(|A|)-(z_A+4)\\
    &\byref{Prop.}{lem:worst-floor-credit}{\ge}3(t_A-1)+(z_A+5)-(z_A+4)
      =3t_A-2\ge4.
\end{align*}
Hence $(B_A+2)/3\ge t_A$; since $t_A$ is an integer,
$d_A=\lfloor(B_A+2)/3\rfloor\ge t_A$.
The padded first block and the remaining comparison blocks, including
the terminal singleton block if needed, satisfy
Proposition~\ref{lem:deadline-cover}. Its ordering has
\[
 g_A\le3d_A-2=3\lfloor(B_A+2)/3\rfloor-2\le B_A.
\]
The first inequality is Proposition~\ref{lem:deadline-cover};
the equality substitutes the computed deadline.
To justify the last inequality, write $B_A+2=3d_A+r_A$ with
$r_A\in\{0,1,2\}$. Then $3d_A-2=B_A-r_A\le B_A$,
so adding $u_A=h_A+e_A^{(4)}$ gives $g_A+u_A\le H^*$.
\end{proof}
The deadline argument bounds every route that uses block scheduling,
including the truncated zero route, by
Proposition~\ref{lem:worst-local-height} and
Claims~\ref{clm:fallback-initial-exits}--\ref{clm:fallback-later-deadline}.
When the hull has dimension at most one, the scheduler is unnecessary:
Lemma~\ref{lem:legal-endpoint-queries} shows that two extreme tests
determine the pattern and leave no endpoints to decode. A nonzero class
then costs at most $2+h_A$. The stored-child, restriction, and
count-credit calculation from Claim~\ref{clm:fallback-initial-exits}
gives
\[
 2+h_A\le3D+\Phi(s)-z_A-3\le H^*.
\]
The second inequality uses $z_A\ge0$ and the definition of $H^*$.
The zero route also ends after at most the two extreme tests. Thus the
local fallback bound holds in every dimension, and induction gives the
stated invariant for the hybrid tree.
\end{proof}
